\documentclass[12pt,a4paper,onecolumn]{article}

% ==============================================================================
% ENCODING, LANGUAGE & GEOMETRY (Cambridge style)
% ==============================================================================
\usepackage[utf8]{inputenc}
\usepackage[british]{babel}

\usepackage{geometry}
\geometry{
  margin=1in,
  headheight=15pt,
  headsep=0.5in,
  footskip=0.5in
}
\usepackage{setspace}
\singlespacing 

% ==============================================================================
% BIBLIOGRAPHY & CITATIONS
% ==============================================================================
\usepackage[backend=biber, style=numeric, sorting=none]{biblatex}
\addbibresource{references.bib}

% ==============================================================================
% GRAPHICS & FIGURES
% ==============================================================================
\usepackage{graphicx}
\graphicspath{{figures/}}
\usepackage{float}
\usepackage{placeins}
\usepackage{subcaption}

% ==============================================================================
% TABLES & ARRAYS
% ==============================================================================
\usepackage{booktabs}
\usepackage{array}
\usepackage{url}
\usepackage{tabularx}
\usepackage{multirow}
\usepackage{threeparttable}
\renewcommand{\arraystretch}{1}
\setlength{\tabcolsep}{8pt}


% ==============================================================================
% MATHEMATICS & TYPOGRAPHY
% ==============================================================================
\usepackage{amsmath}
\usepackage{amssymb}
\usepackage{amsthm}
\usepackage[final]{microtype}
\usepackage{xurl}
\urlstyle{same}

% ==============================================================================
% LISTS & ENUMERATIONS
% ==============================================================================
\usepackage{enumitem}
\setlist[enumerate]{nosep, topsep=0.3em, itemsep=0.3em}
\setlist[itemize]{nosep, topsep=0.3em, itemsep=0.3em}

% HYPERLINKS & METADATA
% ==============================================================================
\usepackage[hidelinks]{hyperref}
\hypersetup{
  pdfborder={0 0 0},
  linktocpage=true,
  colorlinks=false,
  pdftitle={Rugby World Cup 2027 Prediction},
  pdfauthor={Tamara Brand},
  pdfsubject={Applied Data Science}
}

% ==============================================================================
% ENDNOTES & FOOTNOTES
% ==============================================================================
\usepackage{endnotes}
\renewcommand{\notesname}{}
\let\footnote\endnote

% ==============================================================================
% CUSTOM COMMANDS & SHORTCUT CITATIONS
% ==============================================================================
\newcommand{\doi}[1]{\href{https://doi.org}{doi:#1}}

\newcommand{\citeML}{\autocite{Bunker2022}}              % ML in Sports
\newcommand{\citeElo}{\autocite{Elo1978}}                % Elo System
\newcommand{\citeHome}{\autocite{Nevill1999}}            % Home Advantage
\newcommand{\citeGBT}{\autocite{Friedman2001}}           % Gradient Boosting
\newcommand{\citeCal}{\autocite{NiculescuMizil2005}}     % Calibration
\newcommand{\citeCV}{\autocite{Bergmeir2012}}            % Cross-Validation
\newcommand{\citeROC}{\autocite{Fawcett2006}}            % ROC-Analysis

% TITLE & AUTHOR METADATA
\title{Probabilistic Forecasting of International Rugby Match Outcomes:\\
A Machine Learning Approach to the 2027 Rugby World Cup}

\author{Tamara Brand}

\date{\today}

\begin{document}

% --- Title page ---
\begin{titlepage}
  \centering
  {\LARGE\bfseries Probabilistic Forecasting for the 2027 Rugby World Cup:\\[0.3em]}
  {\large\itshape Multi-Source 
  Integration and Uncertainty Quantification\par}
  \vspace{1.6cm}

  {\large
  Tamara Brand\\
  CAS Applied Data Science\\
  University of Bern\\
  \href{mailto:tamara.brand@students.unibe.ch}{tamara.brand@students.unibe.ch}\\[0.8cm]
  \today\par}

  \vspace{1cm}
  \includegraphics[width=0.25\linewidth]{Logo_Springboks.png}
  \vspace{1cm}

\vfill
\begin{abstract}
 
The aim of this project is to develop a reproducible analytical pipeline for predicting the outcomes of international rugby matches and forecasting the 2027
Rugby World Cup (RWC 2027). The analysis integrates three complementary data sources: historical match results from Kaggle \cite{Begbie2022} (Tier 1 nations, 1995–2024); Rugby World Cup (RWC) records from Flashscore \cite{Flashscore2024} (1987–2023); and official World Rugby rankings \cite{WorldRugby2024} (2003–2024).
Feature engineering combines domain knowledge with temporal validity controls \cite{Bergmeir2012} to prevent information leakage. A gradient boosting classifier \cite{Friedman2001} with Elo features \cite{Hvattum2010} achieves 73.4\% accuracy 
and an ROC-AUC of 0.796 \cite{Hanley1982} on held-out test data. Bootstrap confidence intervals \cite{NiculescuMizil2005} quantify uncertainty in model performance.
The primary deliverable is a Monte Carlo simulation of 10,000 tournament paths for the RWC 2027. This yields probabilistic championship and stage
progression estimates for all 24 nations. The results incorporate explicit uncertainty
quantification \cite{Guo2017} and demonstrate how multi-source data integration, temporal validation and transparent modelling can support sports analytics applications 
\cite{Bunker2022}.
 
\end{abstract}
\end{titlepage}

% --- Frontmatter with Roman numerals ---
\pagenumbering{roman}
\setcounter{page}{1}

{\footnotesize
\tableofcontents
}
\clearpage

% --- Main matter with Arabic numerals ---
\clearpage
\pagenumbering{arabic}
\setcounter{page}{1}

\section{Introduction}

\subsection{Background and Motivation}

Over the past three decades, international rugby has become more competitive and data-rich. \cite{OBrien2003} Since the sport's professionalisation in 1995, the organisational landscape of international rugby has undergone substantial transformation. \cite{Hutchins2016} The availability of historical match data and investment in data analytics infrastructure have created significant opportunities for quantitative analysis of match outcomes and tournament dynamics. Predictive modelling in team sports has advanced substantially in recent years. 
\cite{Bunker2022} Research shows that domain-specific features, such as home advantage, recent form, and head-to-head history, improve the accuracy of match predictions. \cite{BunkerThabtah2019,Thomas2008} Elo-based rating systems have been shown to be effective in representing relative team strength in a variety of sports, including rugby union. \cite{Hvattum2010,Elo1978}

Despite these developments, international rugby union remains relatively underexplored compared to sports such as football and basketball. \cite{Bunker2022} Many existing studies rely on limited data sources, static feature sets, or validation approaches that do not adequately account for temporal dependence in match sequences. \cite{Bergmeir2012} Furthermore, only a limited number of studies extend match-level prediction to tournament-level forecasting with explicit uncertainty quantification. \cite{Guo2017,Leitner2010}

Against this backdrop, the focus of this study is South Africa, a country of particular interest given that the team is pursuing an unprecedented third consecutive Rugby World Cup title (following victories in 2019 and 2023; the next tournament is scheduled for 2027). In addition to its analytical relevance, this focus is further motivated by the author's intention to attend the 2027 Rugby World Cup (RWC 2027) in Australia, providing a personal dimension to the investigation of the Springboks' prospects.

\subsection{Research Objectives}
This work addresses these gaps by integrating multiple rugby datasets and Tier 1 match histories from the RWC archive into a unified framework, providing comprehensive coverage of tournaments. Chronological validation is applied to prevent information leakage, and prediction is extended to tournament simulation with transparent uncertainty quantification.

\vspace{0.25cm}

\begin{enumerate}[label=\textbf{RQ\arabic*.}, leftmargin=1.5cm, itemsep=1.5em, parsep=0.25em]
  \item To what extent do domain-specific features, such as Elo ratings, recent form and home advantage, improve the predictive performance of international rugby match outcomes compared to baseline models?
  
  \item Could a forecasting framework that combines machine learning with Elo-based fallback mechanisms produce realistic and well-calibrated probabilistic forecasts for the RWC 2027?
\end{enumerate}

\newpage
\section{Data}

\subsection{Data Sources}

The analysis combines three complementary datasets that together provide information on match outcomes, tournament results and team strength. This integration allows for both match-level feature construction and tournament-specific analysis, reducing reliance on a single source.

\vspace{0.25cm}
\begin{enumerate}[label=\textbf{\arabic*.}, leftmargin=1.5cm, itemsep=0.5em, parsep=0.5em]
  \item \textbf{Kaggle International Rugby Results:} \cite{Begbie2022} covers all matches involving the 10 Tier~1 nations (Argentina, Australia, England, France, Ireland, New Zealand, Scotland, South Africa, Wales, Italy) from 1995 onwards. The professional era baseline (1995) coincides with South Africa's return to the international rugby stage following the end of apartheid-era sporting isolation \cite{WorldRugby1995}.
  
  \item \textbf{Flashscore Rugby World Cup Archive:} \cite{Flashscore2024} comprises web-scraped results from all 10 RWC tournaments (1987–2023), covering all pool and knockout matches for all participating teams. This source is essential for constructing tournament-specific pedigree features and validating model performance across Tier~2 nations not represented in the Kaggle dataset.
  
  \item \textbf{World Rugby Official Rankings:} \cite{WorldRugby2024} provides weekly Men's Rankings snapshots spanning 2003–2024. These rankings implement a rating system analogous to the classical Elo approach, enabling pre-match rating differences to be accurately reconstructed on a match-by-match basis and forming the foundation for team strength measurements throughout the analysis.
\end{enumerate}

\vspace{0.25cm}

Together, these three sources enable complete tournament-level forecasting across all 24 participating nations, combining bilateral match history (Tier~1) with tournament-wide performance metrics.

The raw datasets are not committed to the repository. Instead, reproducible ingestion pipelines download and validate data locally. Cleaned and feature-engineered datasets are generated deterministically by executing the notebooks sequentially.

\subsection{Data Flow and Processing Pipeline}

As illustrated in Figure~\ref{fig:data-flow}, the data engineering pipeline adheres to a three-stage Bronze–Silver–Gold medallion architecture. \cite{Databricks2022}

\vspace{0.25cm}
\begin{enumerate}[label=\textbf{\arabic*.}, leftmargin=1.5cm, itemsep=0.5em, parsep=0.5em]
\item \textbf{Bronze Layer} stores raw data from all source systems in its original form, providing a reproducible foundation for downstream processing.
\item \textbf{Silver Layer} standardizes and cleans the data through deterministic transformation rules. This stage harmonizes schemas, removes inconsistencies, and produces analysis-ready datasets.
\item \textbf{Gold Layer} creates model-ready analytical features. All features are constructed using only information available prior to each match, ensuring a leakage-free predictive modelling framework. \cite{Bergmeir2012}
\end{enumerate}

\begin{figure}[H]
\centering
\includegraphics[width=\linewidth]{DFD.png}
\caption[Data flow and pipeline architecture]{Data flow architecture illustrating the Bronze-Silver-Gold medallion pattern and the flow from raw data ingestion through cleaning to feature engineering.}
\label{fig:data-flow}
\end{figure}
\FloatBarrier

The architecture is implemented through a sequence of Jupyter notebooks, each responsible for a distinct stage of the analytical workflow:

\vspace{0.5cm}
\begin{enumerate}[label=\textbf{\arabic*.}, leftmargin=1.5cm, itemsep=0.5em, parsep=0.5em]
\item \textbf{00\_data\_ingestion:} Downloads historical match results via \texttt{kagglehub}, scrapes RWC match data from Flashscore using Selenium WebDriver, and collects weekly World Rugby ranking snapshots. Raw datasets are stored in the Bronze layer as Parquet files.

\item \textbf{01\_data\_cleaning:} Standardizes variable formats, removes duplicates, and enforces chronological consistency \cite{Bergmeir2012} across records. Filters records to the competitions and periods relevant for analysis. Outputs curated Silver-layer datasets.

\item \textbf{02\_feature\_engineering:} Generates predictive features including Elo rating differentials \cite{Hvattum2010}, rolling performance metrics, and tournament pedigree indicators. All features are calculated using strictly historical information to prevent information leakage. \cite{Bergmeir2012} A chronological train-test split is applied, producing Gold-layer datasets.

\item \textbf{03\_exploratory\_analysis:} Produces descriptive statistics and visualizations to assess data quality, feature distributions, and relationships between predictors and match outcomes.

\item \textbf{04\_model\_training:} Trains and evaluates Logistic Regression and Gradient Boosting models \cite{Friedman2001} with and without Elo-based features. Model performance is evaluated using accuracy, precision, recall, ROC-AUC \cite{Hanley1982}, and Brier score \cite{Brier1950} metrics. The best-performing specification is selected for deployment.

\item \textbf{05\_model\_application:} Applies the selected model to simulate the RWC 2027, executing 10,000 Monte Carlo tournament paths \cite{Leitner2010} and estimating championship and stage-progression probabilities for all participating teams.
\end{enumerate}

\subsection{Data Model}
The analytical unit is a single international rugby match. Each observation includes the match date, the two teams involved, the scores, the venue (home, away or neutral), the competition context and the outcome. The full analysis comprises 3,410 match observations (1,705 unique matches represented from both team perspectives) spanning 1987–2024. For our South African case study (see Section 1.1), we focus on 328 individual matches (1992–2024), of which approximately 73 fall within the held-out test window (2020–2024).
Key variables are summarised in Table~\ref{tab:features}. All features are strictly lagged: rolling statistics use only prior matches, head-to-head records exclude the current match, and Elo ratings are pre-match values. \cite{Bergmeir2012}

\begin{table}[H]
\centering
\caption{Core feature definitions}
\label{tab:features}
\scriptsize
\renewcommand{\arraystretch}{0.9}
\begin{tabularx}{\textwidth}{lllX}
\toprule
\textbf{Feature} & \textbf{Source} & \textbf{Type} & \textbf{Description} \\
\midrule
elo\_diff\_pre & World Rugby + computed & float & Pre-match Elo rating difference \\
home & Multi-source & binary & Venue advantage indicator \\
rolling\_form\_10 & Lagged match history & float & Win rate over last 10 matches \\
rolling\_margin\_3 & Lagged match history & float & Mean score margin (last 3 matches) \\
h2h\_winrate & Lagged match history & float & Historical win rate vs. current opponent \\
consecutive\_wins & Lagged match history & int & Winning streak prior to match \\
rwc\_pedigree & Flashscore RWC data & float & Cumulative RWC performance metric \\
tournament\_tier & Multi-source & categorical & Competition context (RWC, Championship, etc.) \\
\bottomrule
\end{tabularx}
\end{table}

\subsection{Data Quality and Limitations}

The Kaggle and Flashscore datasets are sufficiently structured for match-level analysis, with consistent availability of key variables. Data quality checks focused on type, validity, and completeness.

Dates and scores were converted and matches filtered according to predefined inclusion criteria. Missing values were imputed using domain-informed defaults.

\subsubsection{Dataset Size Clarification}

The Bronze layer contains 1,705 unique matches: 1,396 Tier~1 international Kaggle matches and 309 RWC Flashscore matches. The Silver layer supports team-level modelling by representing each match from both team perspectives (3,410 observations). The Gold layer is divided into a training set (1995--2019, 2,918 observations) and a test set (2020--2024, 492 observations), enforcing strict chronological ordering to prevent information leakage. \cite{Bergmeir2012} The South African dataset contains 328 matches across the full study period, of which approximately 73 fall within the 2020--2024 test window.

\subsubsection{Data Quality Metrics}

Missing values were minimal across the pipeline. World Rugby rankings were successfully matched to 2,533 out of 3,410 observations (74.3\% coverage), with the remaining observations being imputed using domain-informed defaults: 0.5 for win-rate features and 0 for score-margin features. No observations were removed, ensuring the retention of all available data (3,410 records).

\subsubsection{Limitations}

A key limitation is the absence of player-level and injury information. \cite{Williams2016,Hurley2016} While the Kaggle dataset primarily focuses on Tier~1 nations, the RWC archive includes a broader range of emerging nations, which could affect the predictive performance of lower-ranked teams. \cite{Bunker2022} Finally, the analysis assumes that team strength relationships remain stable over time. Structural changes in international rugby and the progressive narrowing of competitive margins among nations may reduce the temporal stability and generalisability of the model. \cite{Bergmeir2012}

\section{Exploratory Data Analysis}

\subsection{Feature Signal and Class Distribution}

Exploratory analysis validates that engineered features exhibit meaningful variation between match outcomes and that the dataset is appropriately balanced for classification. The target variable (Springboks win) shows a class distribution of 60.4\% wins and 39.6\% non-wins (draws or losses) across the 328 South African matches, providing reasonable balance for logistic and tree-based classifiers.

The Elo rating differential shows clear separation by match outcome (Figure~\ref{fig:elo-dist}), with winning matches concentrated at positive Elo differentials and losing matches at negative differentials. This strong visual separation confirms that relative team strength, as represented by Elo-based ratings, is a primary predictor of match outcomes. \cite{Hvattum2010}

\begin{figure}[H]
\centering
\includegraphics[width=0.5\linewidth]{eda_elo_distribution.pdf}
\caption{Elo rating differential distribution by match outcome.}
\label{fig:elo-dist}
\end{figure}

\subsection{Contextual and Performance Predictors}

Home advantage and recent form emerge as strong predictors of match outcomes. South Africa won 75.6\% of home matches ($n=164$) compared with 49.4\% of away matches ($n=168$), representing a difference of 26.2 percentage points. This venue effect is consistent with findings from the sports analytics literature and supports the inclusion of venue-related features in the predictive modelling framework. \cite{Thomas2008, Nevill1999}

Recent performance also exhibits substantial predictive value. Matches preceded by stronger results over the previous three-match window were associated with higher probabilities of subsequent victory, indicating that short-term form captures information beyond longer-term measures of team strength.

\vspace{1em}

Both effects are illustrated in Figure~\ref{fig:eda-contextual}. Home matches show a clear increase in win probability relative to away fixtures, while rolling form demonstrates noticeable separation between winning and non-winning outcomes. Together, these findings provide empirical support for the inclusion of venue and recent-performance features in the predictive models.

\begin{figure}[H]
\centering
\begin{subfigure}{0.48\linewidth}
    \centering
    \includegraphics[width=\linewidth]{eda_home_advantage.pdf}
    \caption{Home advantage effect}
\end{subfigure}
\hfill
\begin{subfigure}{0.48\linewidth}
    \centering
    \includegraphics[width=\linewidth]{eda_rolling_form.pdf}
    \caption{Rolling performance momentum}
\end{subfigure}
\caption[Contextual and performance predictors]{Exploratory analysis of home advantage (left) and rolling performance momentum (right) as predictors of match outcomes.}
\label{fig:eda-contextual}
\end{figure}
\FloatBarrier

\subsection{Elo Rating Dynamics}

The Elo rating system provides a dynamic measure of team strength that evolves with match outcomes. \cite{Elo1978} Matches involving teams with positive Elo differentials (favourable rating advantage) consistently show higher win frequencies, validating the Elo-based approach. \cite{Hvattum2010} The temporal evolution of Elo ratings across the 1995--2024 period reveals distinct competitive cycles, with peaks corresponding to RWC championship years and troughs corresponding to periods of transition or rebuilding. Additional feature analysis is provided in Appendix Figures~\ref{fig:app-pedigree}, \ref{fig:app-correlation}, \ref{fig:app-trend}, and \ref{fig:app-margin}.

\section{Methods}

\subsection{Feature Engineering}

The feature set comprises 13 features organized into three categories reflecting different dimensions of match performance (detailed in Table~\ref{tab:features}). The following sections describe each category.

\vspace{0.25cm}

\begin{enumerate}[label=\textbf{\arabic*.}, leftmargin=1.5cm, itemsep=0.5em, parsep=0.5em]

\item \textbf{Team Strength:} The Elo rating differential serves as the principal measure of relative team strength. \cite{Elo1978} Elo ratings are updated sequentially after each match using only information available at the time of prediction. The expected probability of victory for team A is calculated as:

\begin{equation}
P_A = \frac{1}{1 + 10^{(R_B-R_A)/400}}
\label{eq:elo}
\end{equation}

where $P_A$ denotes the expected probability of victory for team A, and $R_A$ and $R_B$ represent the pre-match Elo ratings. \cite{Elo1978}

\item \textbf{Context and Momentum:} Four variables capture short-term performance and match context. Recent form is measured as rolling win rate over 3, 5, and 10 previous matches. Score momentum is defined as the mean score margin across the previous three matches. Consecutive wins capture winning streak effects, while home advantage is represented by a binary indicator reflecting venue-related performance differences. \cite{Thomas2008,Nevill1999}

\item \textbf{Opponent History and Tournament Pedigree:} Head-to-head win rate captures opponent-specific effects that may not be fully reflected in overall team strength measures. \cite{BunkerThabtah2019} The RWC pedigree metric, derived from tournament records between 1987 and 2023, provides a summary measure of historical RWC performance and long-term tournament success.

\end{enumerate}

\subsection{Model Selection and Validation}

Two base models are fitted: (1) Logistic Regression, chosen for interpretability and transparency; and (2) Gradient Boosting Classifier \cite{Friedman2001}, chosen for its non-linear capacity and reproducible implementation. The GBT model uses the following hyperparameters: 
$n\_\text{estimators}=200$, $\text{learning\_rate}=0.05$, 
$\text{max\_depth}=4$. Each model is trained with and without the Elo feature, yielding four configurations.

\vspace{0.25cm}

To prevent information leakage and respect temporal dependencies, both a chronological train-test split and walk-forward cross-validation are employed: \cite{Bergmeir2012}

\vspace{0.25cm}

\begin{enumerate}[label=\textbf{\arabic*.}, leftmargin=1.5cm, itemsep=0.5em, parsep=0.5em]

\item \textbf{Hold-out Test Set:} Matches from 2020--2024 ($n=492$) are withheld entirely from model training and hyperparameter selection, serving exclusively as a final, unbiased evaluation set.

\vspace{0.25cm}
\item \textbf{Walk-Forward Cross-Validation:} We employ a walk-forward expanding-window approach with four annual test windows (2016--2019). Training data grows progressively from 1995 onwards, reflecting realistic operational scenarios. Each fold extends the training period by one year and evaluates the immediately following season:

\begin{quote}
\small

\textbf{Fold 1:} Training period 1995--2015 ($n=2{,}312$); evaluation on 2016 ($n=112$)

\textbf{Fold 2:} Training period 1995--2016 ($n=2{,}424$); evaluation on 2017 ($n=98$)

\textbf{Fold 3:} Training period 1995--2017 ($n=2{,}522$); evaluation on 2018 ($n=112$)

\textbf{Fold 4:} Training period 1995--2018 ($n=2{,}634$); evaluation on 2019 ($n=156$)
\end{quote}
\end{enumerate}

\vspace{0.25cm}

This dual-layer approach prevents the model from learning on future information while assessing performance stability across multiple temporal windows. The walk-forward results (mean AUC: LR=0.862, GBT=0.860) confirm that the hold-out test performance is representative of model robustness.

\subsection{Model Evaluation}

Performance is assessed using multiple complementary metrics appropriate for classification and probabilistic prediction:

\vspace{0.25cm}

\begin{enumerate}[label=\textbf{\arabic*.}, leftmargin=1.5cm, itemsep=0.25em, parsep=0.5em]

\item \textbf{Accuracy:} measures the overall fraction of correctly classified match outcomes and provides a straightforward benchmark for binary classification performance.

\item \textbf{Precision and Recall:} are reported separately to account for potential class imbalance and to characterise the trade-off between false positives and false negatives.

\item \textbf{ROC-AUC:} serves as a threshold-independent measure of discriminative ability, reflecting the probability that the model correctly ranks a randomly selected positive instance above a negative one. \cite{Hanley1982}

\item \textbf{Brier Score:} constitutes a proper scoring rule for probabilistic forecasts, penalising both overconfident and poorly calibrated predictions; lower values indicate superior probabilistic accuracy. \cite{Brier1950}

\item \textbf{Calibration Curves:} provide a graphical comparison of predicted win probabilities against empirically observed outcome frequencies, enabling visual assessment of systematic over- or underestimation. \cite{Guo2017}
\end{enumerate}

\subsection{Tournament Simulation}

Pre-match win probabilities are derived using a hybrid forecasting strategy. For matchups involving South African opponents present in the training set, the trained GBT model provides team-specific predictions incorporating lagged features and match context; for matchups involving unseen teams or out-of-sample configurations, win probabilities fall back to the Elo logistic function based on official World Rugby rankings. This hybrid approach ensures complete tournament coverage across all 24 teams while leveraging the GBT model's superior discriminative performance (AUC 0.796) where applicable, with the Elo fallback mechanism preventing model failure on unseen teams.

Individual match outcomes are sampled stochastically according to the Bernoulli model: $\text{outcome} \sim \text{Bernoulli}\!\left(P(\text{home win})\right)$, where the success probability is determined by the forecasting strategy described above. The full tournament bracket is simulated across $10{,}000$ independent Monte Carlo replicates, adhering strictly to the official pool structure and knockout draw. The championship probability for each participating nation is estimated as the proportion of simulation replicates in which that team wins the tournament final, while stage-specific progression probabilities (group stage, quarter-final, semi-final, and final) are computed by aggregating outcomes across all $10{,}000$ simulation paths, yielding a complete probabilistic tournament profile for each team.

This approach propagates uncertainty through all knockout rounds and provides empirically grounded probabilistic estimates for all 24 teams. \cite{Leitner2010}

\section{Results}

\subsection{Model Performance}
The chronological test set (2020–2024, $n=492$ matches across all teams) includes matches from the RWC 2023 and 2024 Rugby Championship, periods of competitive volatility that challenge prediction models. This corresponds to approximately 73 South African matches within the held-out period. These results are summarized in Table~\ref{tab:results}.

\begin{table}[H]
\centering
\caption{Model performance on chronological test set}
\label{tab:results}
\footnotesize
\begin{tabularx}{\textwidth}{lcccccc}
\toprule
\textbf{Model} & \textbf{Accuracy} & \textbf{Precision} & \textbf{Recall} & \textbf{ROC-AUC} & \textbf{Brier} \\
\midrule
Logistic Regression & 0.724 & 0.713 & 0.719 & 0.776 & 0.200 \\
Logistic Regression + Elo & 0.724 & 0.709 & 0.727 & 0.789 & 0.193 \\
Gradient Boosting & 0.724 & 0.748 & 0.647 & 0.783 & 0.194 \\
\textbf{Gradient Boosting + Elo} & \textbf{0.734} & \textbf{0.749} & \textbf{0.677} & \textbf{0.796} & \textbf{0.190} \\
\bottomrule
\end{tabularx}
\end{table}

The gradient boosting classifier with Elo features achieves the best overall performance: 73.4\% accuracy, ROC-AUC of 0.796, and Brier score of 0.190. This model is selected as the primary classifier. The Brier score improvement from 0.200 (LR baseline) to 0.190 (GBT+Elo), demonstrating modest gains in probabilistic calibration..

\subsection{Uncertainty Quantification}
The held-out test period (2020--2024, $n=492$ matches across all teams) provides an unbiased evaluation set. For South Africa specifically (n = 73 matches within this window), performance estimates carry additional sampling uncertainty. Bootstrap resampling with replacement (1,000 replicates) is used to estimate 95\% confidence intervals around key metrics. These intervals are displayed in Table~\ref{tab:ci-metrics}.


\begin{table}[H]
\centering
\caption{Best model (GBT+Elo) performance with 95\% bootstrap confidence intervals}
\label{tab:ci-metrics}
\footnotesize
\setlength{\tabcolsep}{12pt}
\renewcommand{\arraystretch}{1.3}
\begin{tabularx}{\textwidth}{lXX}
\toprule
\textbf{Metric} & \textbf{GBT+Elo Estimate} & \textbf{95\% Confidence Interval} \\
\midrule
Accuracy & 0.734 & [0.6951, 0.7724] \\
ROC-AUC  & 0.796 & [0.7557, 0.8322] \\
Brier    & 0.190 & [0.1680, 0.2119] \\
\bottomrule
\end{tabularx}
\end{table}

The wide confidence intervals reflect the inherent unpredictability of international rugby match outcomes. The ROC-AUC interval [0.7557, 0.8322] indicates that while the model reliably discriminates between winners and non-winners, substantial uncertainty remains. This is not a model limitation, but rather reflects genuine competitive volatility in international rugby.
For the RWC 2027 tournament simulation, championship probability estimates similarly carry sampling uncertainty. South Africa's estimated championship probability of 23.8\% is derived from 10,000 Monte Carlo replicates and carries an approximate 95\% confidence interval of [22.4\%, 25.3\%], reflecting both sampling variability and match-level outcome randomness.

\subsection{Feature Importance}
 
The gradient boosting model provides transparent feature importance rankings through permutation-based analysis. Elo rating differential emerges as the dominant predictor (importance: 0.0992), followed by experience differential (0.0625) and home advantage (0.0587). This ranking confirms that relative team strength, contextual venue effects, and historical experience are the primary drivers of match outcomes.

As shown in Figure~\ref{fig:feature_importance}, these results align with domain expectations and justify the focus on team strength and contextual factors rather than more complex nonlinear interactions.
 
\begin{figure}[H]
    \centering
    \includegraphics[width=0.5\linewidth]{model_feature_importance.pdf}
    \caption[Feature importance ranking]{Permutation-based feature importance for the gradient boosting model.}
    \label{fig:feature_importance}
\end{figure}
\FloatBarrier

\subsection{Model Performance: Discrimination and Calibration}

\subsubsection{Discriminative Ability (ROC-AUC)}

To assess the robustness of the selected model, Figure~\ref{fig:roc-ci} displays ROC curves for all four model configurations with 95\% bootstrap confidence intervals (1,000 replicates). The gradient boosting classifier with Elo features (AUC = 0.796, 95\% CI: [+0.0049, +0.0585]) demonstrates superior discrimination compared to the Logistic Regression baseline (AUC = 0.776, 95\% CI: [+0.0049, +0.0585]). The confidence intervals account for both sampling variability in the test set and inherent uncertainty in match outcomes.

\begin{figure}[H]
\centering
\includegraphics[width=0.6\linewidth]{model_roc_curves.pdf}
\caption[ROC curves with bootstrap confidence intervals]{ROC curves for all model configurations, along with their 95\% bootstrap confidence intervals.}
\label{fig:roc-ci}
\end{figure}
\FloatBarrier

\subsubsection{Probabilistic Calibration}

Probabilistic calibration, or the degree of agreement between predicted and observed probabilities, is essential for reliable tournament simulation. The GBT+Elo model demonstrates good calibration across the test set.

\vspace{0.25cm}

\begin{enumerate}[label=\textbf{\arabic*.}, leftmargin=1.5cm, itemsep=0.5em, parsep=0.5em]

\item \textbf{Expected Calibration Error (ECE):} ECE = 0.0808, indicating well-calibrated probability estimates. This metric quantifies the mean absolute difference between predicted and observed win probabilities.\,\autocite{Guo2017}
  
\item \textbf{Brier Score:} 0.190, a proper scoring rule measuring quadratic difference between predicted probabilities and observed outcomes. Lower is better (perfect: 0.0) \,\autocite{Brier1950}.
  
\item \textbf{Calibration Curve:} 
The observed relationship between predicted and actual win probabilities shows slight overconfidence for predictions below 0.6 and slight underestimation above 0.75. Overall agreement is satisfactory for forecasting (Figure~\ref{fig:calibration}).
\end{enumerate}

\vspace{0.25cm}

These metrics collectively indicate that the model's confidence levels 
appropriately reflect true match outcome probabilities, validating its 
use for large-scale tournament simulation where probabilistic accuracy 
is critical.

\begin{figure}[H]
\centering
\includegraphics[width=0.7\linewidth]{model_calibration.pdf}
\caption[Calibration assessment]{Calibration curve (GBT+Elo) showing predicted vs. observed win probabilities. Points near the diagonal indicate good calibration. The dashed line represents perfect calibration.}
\label{fig:calibration}
\end{figure}
\FloatBarrier

\subsection{2027 Rugby World Cup Probabilistic Forecast}

\subsubsection{South African Tournament Progression}

For the case study focusing on South Africa, the stage progression probabilities shown in Figure~\ref{fig:sa-funnel} quantify the likelihood of advancing beyond each knockout phase.

\begin{figure}[H]
\centering
\includegraphics[width=0.5\linewidth]{rwc2027_sa_funnel.pdf}
\caption[SA stage progression probabilities]{South Africa's stage progression probabilities through the RWC 2027 knockout rounds.}
\label{fig:sa-funnel}
\end{figure}
\FloatBarrier

The probabilities reveal a clear tournament gradient: group stage qualification is nearcertain (99.8\%), while quarter-final progression becomes contingent on opponent draw (60.0\%). The semi-final probability drops to 35.3\%, followed by a final stage probability of 12.1\%. The championship estimate of 23.2\% reflects South Africa's competitive but not overwhelming advantage against elite opposition.

\subsubsection{Championship Probabilities Across All Teams}

Figure~\ref{fig:championship} illustrates championship probabilities for all 24 participating teams from the Monte Carlo simulation.

\begin{figure}[H]
\centering
\includegraphics[width=0.6\linewidth]{rwc2027_championship.pdf}
\caption[RWC 2027 championship probabilities]{Championship probabilities for all 24 RWC 2027 teams from 10,000 Monte Carlo simulations.}
\label{fig:championship}
\end{figure}
\FloatBarrier

South Africa receives the highest estimated championship probability (23.2\%) within the simulation framework, followed by New Zealand (16.2\%) and Ireland (15.1\%). These top three nations account for 54.5\% of the simulated championship outcomes, reflecting their consistent high Elo ratings.

\textbf{\textit{Stochasticity Note.}} 
Due to the random sampling of match outcomes in the Monte Carlo procedure, championship probabilities exhibit slight variations ($\pm 0.5\text{--}1.0\%$) across independent simulation runs. The reported values represent a single realization with fixed random seed (\texttt{SEED=42}); confidence intervals from multiple runs would provide more robust uncertainty quantification.

\section{Discussion}

\subsection{Model Performance and Limitations}
The gradient boosting model achieved an accuracy of 73.4\% on the held-out test set, which is a modest but meaningful improvement on the baseline class rate of 62.9\%. With a ROC-AUC of 0.796, the model can be said to reliably discriminate between winners and non-winners. Although Elo ratings are the most important predictor (with an importance score of 0.0992), the non-linear transformation and probabilistic calibration refinements provided by gradient boosting reduce the Brier score from 0.200 to 0.190, demonstrating measurable value in quantifying uncertainty beyond simple rating-based baselines. However, predictions remain probabilistic and individual match outcomes are inherently uncertain.

The moderate Brier score (0.190) reflects both genuine match unpredictability and the challenges of high-stakes international competition. The test period (2020–2024) includes the RWC 2023, which is expected to be one of the most competitive tournaments on record. This naturally constrains the accuracy of predictions.

The absence of player-level, tactical, and injury data represents a fundamental limitation. \cite{Williams2016,Hurley2016} The model captures historical form and contextual factors but not real-time team composition or preparation status. Future improvements should prioritize the integration of squad-level data and tournament-specific contextual variables.

The relatively limited size of the South African case-study dataset introduces additional uncertainty regarding parameter stability and generalisation performance. Although bootstrap confidence intervals partially address this issue, larger datasets covering a broader range of international teams and competitions would likely improve robustness and reduce estimation variance.

\subsection{Tournament Simulation Validity}

The Monte Carlo simulation yields results consistent with expert expectations and recent historical performance. \cite{Leitner2010} South Africa, New Zealand, and Ireland are correctly identified as the top contenders, and the estimated probabilities align reasonably with bookmaker odds and expert consensus forecasts (as of June 2026).

The approach is transparent and auditable: every match probability is derived from official World Rugby rankings using a simple, well-documented Elo formula. \cite{Elo1978} The 10,000 replicates provide sufficient sampling to estimate probabilities to within 1--2 percentage points for top contenders.

\subsection{External Validity and Inherent Limitations}

The predictive framework demonstrates reasonable performance on historical data, but a fundamental limitation merits acknowledgement: sport remains intrinsically unpredictable at the elite level. Several unmeasurable factors shape tournament outcomes.

\subsubsection{Player Injuries and Contingency}

International rugby's collision nature creates inherent injury risk. Key player losses disrupt team composition and psychology simultaneously. Research on the RWC 2015 documented that injury-induced personnel losses increase remaining players' injury fear and reduce confidence. \cite{Hurley2016} Such contingencies remain unmeasurable in historical data yet profoundly influence outcomes.

\subsubsection{Beyond Historical Patterns}

Sports outcomes depend on factors beyond model scope: sleep, nutrition, psychological resilience, and tactical adaptation fluctuate unpredictably. While Elo ratings and recent form capture measurable team strength, future tournaments may differ due to changes in squad composition, coaching strategies, and tactical developments. \cite{Hvattum2010}

\subsubsection{Uncertainty as Information}

The wide confidence intervals (ROC-AUC 95\% CI: [0.756, 0.832]) reflect genuine sporting volatility, not model weakness. South Africa's 23.8\% championship probability is a conditional forecast: given the current information, South Africa is favoured, but in 76.2\% of plausible tournaments, another team would be expected to succeed. These probabilities remain valid only if past outcome drivers remain stable.

\section{Conclusion}

This project involved developing and evaluating a reproducible analytical framework for forecasting international rugby outcomes. By integrating historical match results, \cite{Begbie2022} RWC records \cite{Flashscore2024} and official World Rugby rankings, \cite{WorldRugby2024} the study shows how diverse sports datasets can be transformed into a consistent predictive modelling process while preserving temporal validity and reproducibility.

The findings suggest that measures of relative team strength, contextual variables and recent performance offer valuable predictive insights into international rugby matches. \cite{BunkerThabtah2019} The best-performing model achieved an ROC-AUC score of 0.796 on a chronologically separate test set, and was then used for a large-scale Monte Carlo simulation of the RWC 2027.

Although substantial predictive uncertainty remains, the resulting framework provides a transparent and interpretable approach to tournament forecasting. Rather than attempting to predict exact outcomes, the methodology quantifies uncertainty and supports the probabilistic assessment of alternative tournament scenarios. \cite{Guo2017}

Beyond the RWC application, this study's main contribution lies in its reproducible methodology. Future research could extend the framework by incorporating player-level information, injury reports, squad selection variables, advanced match statistics and continuously updated forecasting pipelines.

The entire workflow, encompassing data acquisition, feature engineering, model training, evaluation and tournament simulation, is entirely reproducible and accessible via the associated GitHub repository.

This framework achieves an accuracy of 73.4\% and an ROC-AUC of 0.796 on historical data, which is a respectable performance for international rugby predictions. However, these metrics only describe the patterns contained in the past. Outside the scope of the model lie unmeasured factors such as injuries to key personnel, tactical innovations, psychological factors and changes to squad composition. These elements continue to be genuine sources of outcome volatility.

The RWC 2027 will unfold in ways that statistical models cannot foresee. Probabilistic forecasting is valuable because it acknowledges uncertainty transparently, quantifies historical patterns reproducibly, and recognises the inherent limits of predicting outcomes for elite competitive events. In that spirit, this framework aims to facilitate structured discussion of scenarios rather than making deterministic forecasts.

\section{Acknowledgements}
I would like to express my gratitude to the CAS Applied Data Science faculty at the University of Bern for their guidance throughout this project. I also thank my study group, ``The Berner Scorpions'', for their constructive exchange and motivation during the development of this work.

\section{Statement}
I declare that I wrote this work independently and used no sources other than those indicated. All parts taken from sources, directly or indirectly, are identified as such. I am aware that otherwise the work will be considered not fulfilled and that the University of Bern may withdraw any degree or title granted on the basis of this work. For verification purposes, I grant the University of Bern the right to process the necessary personal data and to use the work, including reproduction and permanent storage in a database, and to make it available for plagiarism checks.

Artificial intelligence tools were used as methodological and writing support for language refinement, structural feedback, and code documentation. All conceptual decisions, analyses, and interpretations were developed and verified independently by the author.

\newpage
\printbibliography

\appendix
\clearpage
\pagenumbering{roman}
\setcounter{page}{4}

\section{Appendix: Supplementary Figures and Analysis}

\subsection{Exploratory Data Analysis: Supplementary Findings}

The following figures provide additional distributional, temporal, and correlation 
analysis supporting the main EDA findings.

\begin{figure}[H]
\centering
\includegraphics[width=0.7\linewidth]{eda_rwc_pedigree.pdf}
\caption[RWC pedigree and win rate relationship]{Relationship between cumulative RWC performance and overall win rate.}
\label{fig:app-pedigree}
\end{figure}

\begin{figure}[H]
\centering
\includegraphics[width=0.7\linewidth]{eda_correlation_matrix.pdf}
\caption[Feature correlation matrix]{Correlation structure among engineered features, revealing moderate correlations between Elo-based and rolling performance metrics.}
\label{fig:app-correlation}
\end{figure}

\begin{figure}[H]
\centering
\includegraphics[width=0.7\linewidth]{eda_score_margin.pdf}
\caption[Score margin distribution]{Score margin distribution for South Africa by match outcome.}
\label{fig:app-margin}
\end{figure}

\begin{figure}[H]
\centering
\includegraphics[width=0.7\linewidth]{eda_winrate_trend.pdf}
\caption[Annual win rate trend]{South African annual win rates from 1995 to 2024 with 3-year smoothing.}
\label{fig:app-trend}
\end{figure}

\subsection{Tournament Forecast Sensitivity and Details}

\begin{figure}[H]
\centering
\includegraphics[width=0.65\linewidth]{rwc2027_sa_path.pdf}
\caption[Tournament bracket sensitivity]{South Africa's championship probability under different knockout draw scenarios.}
\label{fig:app-path}
\end{figure}

\begin{figure}[H]
\centering
\includegraphics[width=0.65\linewidth]{rwc2027_tier1.pdf}
\caption[Tier 1 championship comparison]{Championship probabilities for the 10 Tier~1 nations.}
\label{fig:app-tier1}
\end{figure}

\end{document}